% !TeX root = ../../Book2.tex
\section{最大解、延拓与连续依赖}
\label{b2:sec:ode-continuation}

短时间构造不规定解的最终时间范围。唯一性使不同局部构造能够相互接合；增长估计则用于判断这种接合能否经过任意有限时间。本节先在自治 \(C^1\) 条件下整理这两个步骤。

\subsection{最大解区间}
\label{b2:subsec:ode-maximal-interval}

设 \(U\subset\mathbb R^d\) 开，\(b\in C^1(U;\mathbb R^d)\)、\(a\in U\)。考虑所有在含 \(0\) 的开区间上满足 \(X'=b(X)\)、\(X(0)=a\) 的解。它们在重叠区间上由\cref{b2:thm:07-local-ode}的唯一性相等，因此可以拼成一条解。所有这些区间的并仍为含 \(0\) 的开区间，记为 \(I_a\)。每个时刻附近的拼接函数都等于某条局部解，所以它属于 \(C^1\) 并满足方程。这就得到经过 \(a\) 的最大解；任何同初值解都是它的限制。

\subsection{共同短区间上的初始点依赖}
\label{b2:subsec:ode-initial-dependence}

在\cref{b2:thm:07-local-ode}选定的共同紧球和共同时间 \([-h,h]\) 上，式\eqref{b2:eq:07-flow-lipschitz}给出
\[
 \sup_{|t|\leq h}|X(t,a)-X(t,\widetilde a)|
 \leq e^{Lh}\cdot|a-\widetilde a|.
\]
因此解在一致范数下连续地依赖初始点，且此处具有 Lipschitz 控制。这个共同时间来自局部球上的速度界和导数界，不能直接当作所有初始点的共同全局时间。

\subsection{线性增长与全时间存在}
\label{b2:subsec:ode-linear-growth}

若速度至多随位置线性增长，Gronwall 估计可以把任意有限时间内的轨道限制在紧球中。

\begin{proposition}\label{b2:prop:ode-linear-growth}
设 \(d\) 为正整数，\(b\in C^1(\mathbb R^d;\mathbb R^d)\)，且存在 \(C\geq0\) 使对所有 \(x\in\mathbb R^d\) 都有
\[
 |b(x)|\leq C(1+|x|).
\]
则对每个 \(a\in\mathbb R^d\)，初值问题 \(X'=b(X)\)、\(X(0)=a\) 的最大解在整个实轴上存在，并满足
\[
 1+|X(t,a)|\leq(1+|a|)\cdot e^{C|t|}.
\]
\end{proposition}
\begin{proof}
对正时间，积分方程给出
\[
 1+|X(t,a)|\leq1+|a|+C\int_0^t(1+|X(s,a)|)\,\mathrm ds.
\]
在最大解区间内的任意紧时间区间上应用\cref{b2:lem:ode-gronwall}，得到所述估计；负时间对反向速度 \(-b\) 使用同一论证。若最大解区间有有限右端点 \(T\)，此估计使轨道在 \([0,T)\) 落于某个紧球。该球上速度有界，所以 \(|X(t,a)-X(s,a)|\leq M|t-s|\)，从而 \(t\uparrow T\) 时有极限 \(a_T\)。由 \(a_T\) 出发的局部解可与原解拼接；积分方程保证拼接点导数也满足方程，与最大性矛盾。左端点相同。
\end{proof}

线性增长条件不能无条件省去。对于 \(b(x)=x^2\)，从 \(a>0\) 出发的轨道为 \(a/(1-at)\)，在 \(t=1/a\) 逃向无穷。局部唯一性仍成立，但这条最大解不能延续到所有正时间。另一方面，若只关心一条有限轨道及其附近，紧集上的局部条件已经足够，不必强求整个空间上的有界导数。

\subsection{一般延拓准则}
\label{b2:subsec:ode-general-continuation}

待撰写。拟在一般定义域中区分离开每个紧集、到达定义域边界和状态无界等情形，并补充非自治问题的延拓与连续依赖。

\subsection{标量比较}
\label{b2:subsec:070207}

待撰写。
